\sect{Ćwiczenia 10}

\begin{frame}{Zadanie 1}
	\begin{figure}
		\includegraphics[scale=0.08, alt={}]{./imags/E10/01.png}
		\caption{Zadanie 1. Oblicz napięcie na kondensatorze.}
	\end{figure}	
\end{frame}

\begin{frame}{Zadanie 2}
	\begin{figure}
		\includegraphics[scale=0.08, alt={}]{./imags/E10/02.png}
		\caption{Zadanie 2. Oblicz prąd cewki.}
	\end{figure}	
\end{frame}

\begin{frame}{Zadanie 3}
	\begin{figure}
		\includegraphics[scale=0.08, alt={}]{./imags/E10/03.png}
		\caption{Zadanie 3. Oblicz prąd cewki, napięcie na kondensatorze i częstotliwość rezonansową obwodu.}
	\end{figure}	
\end{frame}

\begin{frame}{Zadanie 4}
	\begin{figure}
		\includegraphics[scale=0.08, alt={}]{./imags/E10/04.png}
		\caption{Zadanie 4. Oblicz prąd cewki, napięcie na kondensatorze i częstotliwość rezonansową obwodu.}
	\end{figure}	
\end{frame}

\begin{frame}{Zadanie 5}
	\begin{figure}
		\includegraphics[scale=0.08, alt={}]{./imags/E10/05.png}
		\caption{Zadanie 5. Wykorzystując impedancję operatorową oblicz napięcie na kondensatorze i prąd cewki.}
	\end{figure}	
\end{frame}
